\documentclass[11pt]{article}
\usepackage[a4paper,margin=1in]{geometry}
\usepackage{amsmath,amssymb,amsthm,hyperref}
\newtheorem{theorem}{Theorem}
\newtheorem{corollary}{Corollary}
\newtheorem{lemma}{Lemma}
\title{GST cyclic causal folds acting on the original variety:
an all-arity Hodge-genesis and purity-defect law}
\date{22 October 2026}
\begin{document}
\maketitle

\begin{abstract}
This is a substantive extension of
\texttt{GST\_causal\_graph\_double\_fold\_weight\_emergence\_2026-10-22.tex}.
Instead of adjoining a singular causal curve as an independent spectator,
the fold carries a nontrivial algebraic action on cohomology derived from
\emph{every} original connected smooth projective complex variety $X$.
For all fold arities $m\geq2$ we construct a concrete projective
twisted-nodal realization and a smooth projective companion; derive exact
mixed/pure Hodge-generating equations; and prove that two causal loops
produce an even-degree, rational diagonal but non-pure sector depending on
the original $X$.  The original $X$ is input geometry, not an assumed
cycle-class solution.  These equations do not assert a Hodge counterexample
on a smooth projective variety.
\end{abstract}

\section{Immutable original-$X$ contract and dimensionless input}
Fix any \emph{connected}, smooth, projective complex variety $X$
of complex dimension $d\geq0$.  The immutable rational Hodge target
would require that for every $p$ and rational $(p,p)$ Hodge class
$\alpha$ one constructs $Z_\alpha\in CH^p(X)_{\mathbb Q}$
with $\mathrm{cl}^{\mathrm{can}}_X(Z_\alpha)=\alpha$.
We make no algebraicity assumption about $\alpha$.

The GST input is instead a finite causal incidence cycle, its cyclic
automorphism and a geometric realization.  World arity $m$ is chosen
freely; geometry then produces the cohomological degrees and Hodge types.
In the language of \texttt{GSTGeneralSpace.lean}, the underlying world
is not assigned a metric or a Hodge decomposition.  The construction
below is the explicit algebraic realization required for its
dimension-fold laws to have mathematical content.

For a mixed Hodge structure write
\[
 P_Y(t,s,u,v)=
 \sum_{k,w,p,q}
 h^{p,q}\bigl(\operatorname{Gr}^W_wH^k(Y,\mathbb Q)\bigr)
 t^k s^w u^p v^q,\qquad w=p+q.
\]
For smooth projective $Y$, every degree-$k$ sector is pure of
weight $k$.  Define the two GST fold readouts
$\chi=p-q$ and $\delta=w-k$; the former tests diagonality, the latter
the cohomological purity condition.

\section{Actual cyclic original-$X$ causal realization}
Let $m\geq2$, $B=X^m$, and $\rho:B\to B$ cyclically permute factors.
Write $G=C_m=\langle\rho\rangle$.
Let $D_m$ be the connected \emph{cycle of $m$ copies of
$\mathbb P^1$}, with $\infty$ on component $i$ glued transversely
to $0$ on component $i+1$ (indices modulo $m$).
For $m=2$ there are two distinct nodes and two components.
The cyclic rotation $\tau$ of components preserves the nodes;
no nonidentity rotation fixes any geometric point.

Take the actual projective geometric quotient
\[
       Y_m(X)=(X^m\times D_m)/C_m ,
       \qquad
       g\cdot(b,c)=(\rho(g)b,\tau(g)c).
       \tag{1}
\]
The group action is free because the rotation on $D_m$ is free,
even when $\rho$ has fixed points on $B$.
The quotient is projective (finite quotient of a projective scheme)
and has complex dimension $md+1$. It is a singular nodal
normal-crossings realization with a single rational normalization
component $B\times\mathbb P^1$; equivalently its two boundary
copies of $B$ are glued by $\rho$.

For an independent \emph{smooth outcome with the same causal
group and the same action on $X$}, choose a smooth elliptic curve
$E$ and an $m$-torsion point $a$ of exact order $m$.
Translation $\tau_a$ acts freely on $E$. Define
\[
       S_m(X)=(X^m\times E)/C_m ,
       \qquad
       g\cdot(b,e)=(\rho(g)b,e+g a).
       \tag{2}
\]
This quotient is smooth projective of the same complex dimension
$md+1$ and has a pure Hodge structure.
Equations (1) and (2) are two \emph{actual} realized world outputs.
An equivariant flat family interpolating them is a separate
geometric assertion and is not assumed here.

\begin{lemma}[Exact rational cohomology of the causal cycle]
The $C_m$-equivariant mixed Hodge structures on $D_m$ are
\[
 H^0(D_m,\mathbb Q)=\mathbb Q(0),\qquad
 H^1(D_m,\mathbb Q)=\mathbb Q(0),\qquad
 H^2(D_m,\mathbb Q)=\mathbb Q(-1)\otimes\mathbb Q[G].
 \tag{3}
\]
$G$ acts trivially on degrees zero and one, and by the regular
component-permutation representation on degree two.
\end{lemma}
\begin{proof}
The normalization is the disjoint union of $m$ rational
components. The node--component incidence graph is one
oriented $m$-cycle. Its boundary matrix has rank $m-1$;
there is one graph $H^1$, of weight zero and type $(0,0)$.
Rotation preserves its orientation, so acts trivially on that line.
Top cohomology has one weight-two Tate class for each rational
component, and $G$ rotates these $m$ classes regularly.
The normalization sequence is equivariant and compatible with
mixed Hodge structures.
\end{proof}

\begin{theorem}[GST original-$X$ cyclic causal fold decomposition]
Put $V_k=H^k(B,\mathbb Q)$ and $V_k^G=\ker(\rho^*-\mathrm{id})$.
For every integer $k$, with $V_j=0$ for $j<0$, the following is
an isomorphism of rational mixed Hodge structures:
\[
 \boxed{\quad
 H^k(Y_m(X),\mathbb Q)\cong
 V_k^G\ \oplus\ V_{k-1}^G\
                  \oplus\ V_{k-2}(-1).
 \quad}                                                \tag{4}
\]
Here the middle summand occurs in cohomological degree $k$
but has weight $k-1$; the other two summands have weight $k$.
The corresponding smooth companion satisfies
\[
 \boxed{\quad
 H^k(S_m(X),\mathbb Q)\cong
 V_k^G \oplus
 (V_{k-1}^G\otimes H^1(E,\mathbb Q))
 \oplus V_{k-2}^G(-1).
 \quad}                                                \tag{5}
\]
Its three summands all have weight $k$.
\end{theorem}
\begin{proof}
For the free finite cyclic action, rational transfer makes
cohomology of each quotient the $G$-invariants of the respective
product cohomology, compatibly with its mixed Hodge structure.
K\"unneth and (3) give the first two summands in (4). For the third,
the regular representation identity
$(V\otimes\mathbb Q[G])^G\cong V$ holds by the explicit map
$v\mapsto\sum_{g\in G}g v\otimes e_g$.
Its inverse is coefficient extraction at $e_1$, so no
extra invariant-only assumption is inserted.
For (5), translation on an elliptic curve acts identically on
its cohomology. The elliptic cohomology is $\mathbb Q(0)$,
$H^1(E)$, and $\mathbb Q(-1)$ in degrees $0,1,2$.
Finite invariants and K\"unneth now give (5).
\end{proof}

\section{ALL-arity equation, graded Koszul signs, and exact output comparison}
Let $P_X=P_X(t,s,u,v)$ and $A=t^2s^2uv$. For
$0\leq j<m$ put $g_j=\gcd(m,j)$, $L_j=m/g_j$.
The graded K\"unneth action on $H^*(X)^{\otimes m}$ is a
\emph{signed} permutation, with Koszul sign. Consequently
\[
 \boxed{\quad
 I_m(X)=P_{H^*(X^m)^{C_m}}
 =\frac1m\sum_{j=0}^{m-1}
 \left[
   P_X\bigl((-1)^{L_j-1}t^{L_j},
            s^{L_j},u^{L_j},v^{L_j}\bigr)
 \right]^{g_j}.
 \quad}                                                \tag{6}
\]
The $j=0$ term is $P_X^m$.
Indeed $\rho^j$ consists of $g_j$ cycles of length $L_j$;
only tensor-basis words constant around each cycle contribute
to its trace. On a homogeneous degree-$k$ basis element, the
cycle carries sign $(-1)^{(L_j-1)k}$. Averaging the characters
is the rational invariant projector.

The complete all-arity GST dimensional-fold genesis equations are
\[
 \boxed{\begin{aligned}
 P_{Y_m(X)}&=(1+t)\,I_m(X)+A\,P_X^m,\\
 P_{S_m(X)}&=(1+t\,s(u+v)+A)\,I_m(X),\\
 P_{S_m(X)}-P_{Y_m(X)}
     &=t\bigl(s(u+v)-1\bigr)I_m(X)
       +A\bigl(I_m(X)-P_X^m\bigr).
 \end{aligned}}                                       \tag{7}
\]
These are \emph{entire Hodge-graded realizations}, not
endomorphisms relabeling a fixed $H^{p,q}(X)$.
For any additional Tate dimension-fold $N\geq0$, multiply
both series by $\sum_{\ell=0}^N A^\ell$.
Thus arbitrary $m,N$ yield actual dimensions $md+1+N$,
while preserving the exact charge and purity-gap diagnostics.

\section{Genuine original-product Chow operator}
The action $\rho$ is an algebraic automorphism of
$B=X^m$. It supplies actual graph correspondences
$[\Gamma_{\rho^j}]\in CH^{md}(B\times B)_{\mathbb Q}$.
The averaged correspondence
\[
 \boxed{\qquad
 e_\rho=\frac1m\sum_{j=0}^{m-1}[\Gamma_{\rho^j}],
 \qquad e_\rho\circ e_\rho=e_\rho
 \quad\text{in } CH^{md}(B\times B)_{\mathbb Q}.
 \qquad}                                               \tag{8}
\]
is a \emph{genuine Chow idempotent}, because graph composition
implements the cyclic group multiplication. It induces the
rational Hodge projector onto $V_k^G$. Therefore the invariant
factor in the fold formula (7) has an actual source in scheme
correspondence geometry, not an assumed Hodge operator.
Its being algebraic as a \emph{two-legged correspondence}
does not assert an algebraic representative for every
one-legged Hodge class on the untouched original $X$.

\section{Second fold: even-degree diagonal weight failure for EVERY X}
Let $D_1$ denote the irreducible rational nodal projective curve,
so $P_{D_1}=1+t+A$. Construct
\[
             Z_m(X)=Y_m(X)\times D_1.
\]
It is an actual singular projective world of dimension $md+2$,
and
\[
 \boxed{\qquad
 P_{Z_m(X)}=(1+t+A)
             \bigl((1+t)I_m(X)+A P_X^m\bigr).
 \qquad}                                               \tag{9}
\]
The sole \emph{double-weight-deficit} contribution is the
$t^2 I_m(X)$ term. Hence for every $k$
\[
 \boxed{\quad
 \operatorname{Gr}^W_{k}H^{k+2}(Z_m(X),\mathbb Q)
 \cong H^k(X^m,\mathbb Q)^{C_m},
 \quad(\chi,\delta)=(p-q,-2)\text{ on a }(p,q)\text{ component}.
 \quad}                                                \tag{10}
\]
For $d>0$ choose the nonzero rational top point class
$\eta_X\in H^{2d}(X,\mathbb Q)$, of type $(d,d)$.
Inside $H^{2d}(X^m,\mathbb Q)$ the vector
\[
       \Theta_X=\sum_{i=1}^m
       1\otimes\cdots\otimes\eta_X\otimes\cdots\otimes1
                                                               \tag{11}
\]
is nonzero and cyclically invariant. The indices are
even-degree, so no Koszul signs arise.
Equation (10) constructs from \emph{every} original $X$
a rational cohomological sector on $Z_m(X)$ with
\[
       k_{\mathrm{new}}=2d+2,\quad
       (p,q)=(d,d),\quad
       w=2d,\quad
       (\chi,\delta)=(0,-2).
                                                               \tag{12}
\]
After any Tate fold $j$ its signature becomes
$(2d+2+2j,d+j,d+j,2d+2j)$, with the same charge $0$
and defect $-2$. Therefore \emph{charge-zero does not imply
smooth-projective purity}, even with an original-$X$-dependent
algebraic cyclic action and genuine projective realization.
The singular world still has its legitimate \emph{mixed}
Hodge structure; this is not a nonalgebraic rational $(p,p)$
class on a smooth projective $X$.

\section{Exact test: X an elliptic curve, m=2}
For an elliptic curve $X$, $P_X=1+ts(u+v)+A$.
The cyclic involution swaps the two $X$ factors. Formula (6)
specializes to
\[
 I_2=\frac12\bigl(P_X^2+
       P_X(-t^2,s^2,u^2,v^2)\bigr).
\]
The symbolic rational Wolfram evaluation of (7) gives
\[
 \begin{array}{c|c|c}
  &\operatorname{Gr}^W_1 H^2&
        \operatorname{Gr}^W_2 H^2\\ \hline
 Y_2(X)&u+v&3uv\\
 X^2\times D_1&2u+2v&u^2+5uv+v^2\\
 S_2(X)&0&u^2+5uv+v^2
 \end{array}
\]
where displayed polynomials are the $u,v$ Hodge multiplicities.
For the second actual nodal fold, the weight-two component of
$H^4(Z_2(X))$ is $2uv$ and has $\delta=-2$.
A separate check with $X=\mathbb P^1$ gives $uv$ in this
degree-four weight-two component. Thus the signature
is not an artifact of elliptic odd cohomology.

\section{Target-sensitive lift of EVERY original-$X$ Hodge class}
The preceding fold is not limited to a selected point or top
cohomology class. Fix \emph{any} rational Hodge class
$\alpha\in H^{2p}(X,\mathbb Q)\cap H^{p,p}(X)$.
The $G$-invariant cyclic symmetrization
\[
 \Theta_m(\alpha)=\sum_{i=1}^m\operatorname{pr}_i^*\alpha
 \in H^{2p}(B,\mathbb Q)^G\cap H^{p,p}(B)
 \tag{13}
\]
is a rational Hodge class without selecting any source cycle.
For $p>0$, fixing a complex point $x_0\in X(\mathbb C)$
gives an algebraic insertion
$s:X\hookrightarrow B$, $x\mapsto(x,x_0,\ldots,x_0)$, and
\[
 \boxed{s^*\Theta_m(\alpha)=\alpha.}                     \tag{14}
\]
Indeed for $i\ne1$ the composition $\operatorname{pr}_i s$
is constant and pulls every positive-degree class to zero.
For $p=0$, the map sends $\alpha$ to $m\alpha$ and is
injective over $\mathbb Q$. Therefore $\Theta_m$ is injective
for every codimension, with an explicit left inverse.

Combining (13) with the weight-$2p$ K\"unneth summand
in (10) constructs a \emph{linear injection}
\[
 \boxed{\quad
 \mathfrak B_{m,p}:
 H^{2p}(X,\mathbb Q)\cap H^{p,p}(X)
 \hookrightarrow
 \operatorname{Gr}_{2p}^W H^{2p+2}(Z_m(X),\mathbb Q).
 \quad}                                                \tag{15}
\]
Every nonzero target is sent to a nonzero rational diagonal
class in degree $2p+2$, of weight $2p$ and therefore
of purity gap $-2$.  This is a \emph{target-sensitive
nonpure Hodge-sector construction for every original $X$}
that does not assume $\alpha$ is algebraic.  The target
class is not thereby made into a nonalgebraic class of
a smooth projective variety: the receiving world is singular.

\begin{theorem}[Algebraicity-faithful smooth causal fold]
Let $q:B\times E\to S_m(X)$ be the finite \'etale quotient.
There is a unique rational Hodge class $\widehat\alpha$
on the \emph{smooth} $S_m(X)$ such that
\[
 q^*\widehat\alpha=\operatorname{pr}_B^*\Theta_m(\alpha).
 \tag{16}
\]
For every $X,m,p,\alpha$ as above,
\[
 \boxed{\quad
 \alpha\in\operatorname{im}\mathrm{cl}^{\mathrm{can}}_X
 \ \Longleftrightarrow\
 \widehat\alpha\in\operatorname{im}\mathrm{cl}^{\mathrm{can}}_{S_m(X)}.
 \quad}                                                \tag{17}
\]
Both implications are induced by genuine algebraic
pushforward/pullback operations, not a freely chosen
Hodge-to-Chow inverse.
\end{theorem}
\begin{proof}
Finite-group rational transfer identifies $H^*(S_m,\mathbb Q)$
with the invariants in $H^*(B\times E,\mathbb Q)$.
The right side of (16) is invariant under cyclic permutation
of $B$ and torsion translation on $E$; hence it descends
uniquely to $\widehat\alpha$.  All these maps preserve
Hodge type, so $\widehat\alpha$ is rational of type $(p,p)$.

Suppose $\alpha=\mathrm{cl}^{\mathrm{can}}_X(W)$ for
$W\in CH^p(X)_{\mathbb Q}$. Set
$W_B=\sum_i\operatorname{pr}_i^*W\in CH^p(B)_{\mathbb Q}$.
It is invariant as a Chow class. Then the honest rational
algebraic cycle
\[
       \widehat W={1\over m}q_*
            (\operatorname{pr}_B^*W_B)
       \in CH^p(S_m)_{\mathbb Q}
                                                               \tag{18}
\]
has canonical Betti class $\widehat\alpha$:
the finite-degree transfer identity and invariance give
$q^*\mathrm{cl}_{S_m}^{\mathrm{can}}(\widehat W)
=\mathrm{cl}_{B\times E}^{\mathrm{can}}
 (\operatorname{pr}_B^*W_B)$, and $q^*$ is injective.

Conversely suppose $\widehat\alpha
=\mathrm{cl}_{S_m}^{\mathrm{can}}(\widehat W)$ for a
genuine codimension-$p$ cycle $\widehat W$.
Pull it back by the finite \'etale map $q$, then by the
regular embedding $i_e:B\to B\times E$ at any
point $e\in E(\mathbb C)$, obtaining
$W_B=i_e^*q^*\widehat W\in CH^p(B)_{\mathbb Q}$.
Naturality of the canonical cycle class and (16) give
$\mathrm{cl}_B(W_B)=\Theta_m(\alpha)$.
For $p>0$ set $W=s^*W_B$ using the algebraic
insertion $s:X\to B$; then (14) proves
$\mathrm{cl}_X(W)=\alpha$.
For $p=0$, pullback returns $m\alpha$ and
$W=m^{-1}s^*W_B$ gives $\alpha$.
All these operations lie in actual rational Chow groups.
\end{proof}

Equation (17) is a sharp \emph{algebraicity-preservation
equivalence}, not a proof of universal source existence:
an arbitrary preimage $\widehat W$ on $S_m$ is not
constructed from $\alpha$ by (16).  Its value is that
the smooth causal fold \emph{faithfully transports any
unknown original-X Hodge target} and has an explicit
native cycle return whenever a folded native source
is independently available. The singular double fold
(15) instead exposes a nonpure topology-dependent
version of the same unknown class.

\section{Evidence and precise limits}
Uniform arguments: finite-group transfer; rational K\"unneth;
node-normalization mixed Hodge structure; character projection
including Koszul signs; algebraic graph composition.
Finite exact Wolfram checks at $m=1,\ldots,5$ for elliptic
cohomology, at $m=2$ for elliptic and $\mathbb P^1$ double
folds, and polynomial smooth--singular comparison, returned
exact equalities and nonnegative coefficients.
No approximation, unproved abstract cycle-class map,
source preimage, local Lean installation, or claim of a
new universal Hodge solution is involved.

\emph{Mathematical boundary.}
The construction acts nontrivially on the entire Hodge
cohomology of $X^m$, for arbitrary input $X$, and supplies
a genuine Chow projector there.  It changes realized
global topology and realizes both pure and mixed outputs.
It does not yet construct $Z_\alpha\in CH^p(X)_{\mathbb Q}$
for a prescribed arbitrary Hodge class $\alpha$ on the
\emph{same} $X$. The decisive additional operation would
need a genuine Chow correspondence returning an independently
created one-legged class to the untouched original $X$;
this cannot be deduced just from the cyclic invariants.
\end{document}
